\cotefile{43}
\scanpage{79}
% archive sheet number: 78
% L1
\uline{\unsure{Supposons} $\dim A\geq 2$}, \unsure{on aura donc} $H^i_{\unsure{A}}(\unsure{G|\mathcal O})=0$ si $i=0,1$,

% L2--3
\deleted{Alors} $H^2(X/X-o,G)$ \unsure{est une extension de} $\operatorname{Coker}\bigl(\unsure{H^1(X,G)}\to H^1(X-o,G)\bigr)$ \unsure{par} $\operatorname{Ker}\bigl(H^2(X)\to H^2(X-o,G)\bigr)$.

% L4--7; insertion between L4 and L5 points after « propre ».
Si $G$ \unsure{est extension d'un discret par un propre}, \inserted{\unsure{et est} \uline{\unsure{simple|lisse}} \unsure{sur $k$}}, \unsure{alors par le théorème} \uline{de pureté}
\[
H^1(X,G)\to H^1(X-o,G)\quad\text{\unsure{est} \uline{surjectif}}
\]
[et si $G$ \unsure{est radiciel} ? \unsure{\ldots}]. \unsure{Quand|Quant} à $H^2(X)\to H^2(X-o,G)$
